\chapter{LCOV Format Reference}
\label{appx:lcov}

This appendix documents the LCOV \texttt{.info} file format directives that the converter processes, and the \texttt{coverage.json} output schema.

\section{LCOV Directives}

The LCOV format is a line-oriented text format. Each source file section begins with \texttt{SF:} and ends with \texttt{end\_of\_record}. The converter processes the following directives:

\begin{lstlisting}[caption={LCOV .info file format example}]
TN:                          # Test name (ignored)
SF:/path/to/source.c         # Source file path
FN:51,55,function_name       # Function: start,end,name
FNDA:11,function_name        # Function data: count,name
FNF:4                        # Functions found (total)
FNH:3                        # Functions hit (executed)
DA:51,11                     # Line 51 executed 11 times
DA:96,0                      # Line 96 never executed
LF:63                        # Lines found (ignored)
LH:38                        # Lines hit (ignored)
BRDA:51,0,0,1                # Branch data (not yet supported)
BRF:10                       # Branches found (not yet supported)
BRH:7                        # Branches hit (not yet supported)
end_of_record                # End of file section
\end{lstlisting}

The converter uses \texttt{DA:} records as the authoritative source for line executability. Lines with \texttt{DA:N,count>0} are covered; lines with \texttt{DA:N,0} are uncovered but executable. Lines without any \texttt{DA:} record are non-executable.

\section{coverage.json Schema}

\begin{lstlisting}[caption={coverage.json output schema}]
{
  "<absolute_file_path>": {
    "lines": [<line_numbers with count > 0>],
    "uncovered": [<line_numbers with count == 0>],
    "functions_found": <integer from FNF>,
    "functions_hit": <integer from FNH>
  }
}
\end{lstlisting}

The legacy format (plain list of covered line numbers per file) is also supported for backward compatibility:

\begin{lstlisting}[caption={Legacy coverage.json format}]
{
  "<absolute_file_path>": [3, 4, 6, 8, 10]
}
\end{lstlisting}

\section{Database Encoding}

In the \texttt{test\_coverage} table, the \texttt{covered\_lines} column uses sign encoding:

\begin{itemize}
    \item Positive values represent covered lines (e.g., \texttt{51} means line 51 was executed)
    \item Negative values represent uncovered executable lines (e.g., \texttt{-96} means line 96 is executable but was never executed)
    \item Lines not present in the table are non-executable
\end{itemize}

The \texttt{getTestCoverage} API returns all values for a file. The frontend checks the sign to determine coloring: positive $\rightarrow$ green, negative $\rightarrow$ red, absent $\rightarrow$ no color.
